% polyglossia: Russian with CM Unicode (TeX Live's cm-unicode)
\documentclass{article}
\usepackage{fontspec}
\setmainfont{cmunrm.otf}[BoldFont=cmunbx.otf, ItalicFont=cmunti.otf, BoldItalicFont=cmunbi.otf]
\setsansfont{cmunss.otf}[BoldFont=cmunsx.otf, ItalicFont=cmunsi.otf]
\setmonofont{cmuntt.otf}
\usepackage{polyglossia}
\setmainlanguage{russian}
\setotherlanguage{english}
\begin{document}
\section{Введение}
Съешь же ещё этих мягких французских булок, да выпей чаю. Широкая электрификация
южных губерний даст мощный толчок подъёму сельского хозяйства.
\textbf{Полужирный}, \textit{курсив}, \textsf{рубленый}, \texttt{моноширинный}.
\subsection{Формулы}
Если $x>0$, то $\sqrt{x}^2 = x$. \textenglish{English text inside Russian.}
Дата: \today.
\end{document}
